\section{Details for the proof of Theorem~\ref{thm:zero-free}}\label{app:zero-free,details}
\subsection{Preliminaries}
%
Lemma~\ref{lem:kcol} directly implies Theorem~\ref{thm:zero-free}, and we give the proof in this section.
%
%
First, we state a lemma of Barvinok which is useful for evaluating sums of restricted partition functions.
\begin{lemma}[{Barvinok~\cite[Lemma 3.6.3]{Barvinok16}}]
 \label{lem:cone}
 Let $u_1, \ldots, u_n \in \RR^2$ be non-zero vectors such that the angle between any two vectors $u_i$ and $u_j$ is at most $\alpha$ for some $\alpha \in [0,2\pi/3)$.
 Then the $u_i$ all lie in a cone of angle at most $\alpha$ and
 \[
 \bigg| \sum_{i=1}^n u_i \bigg| \geq \cos(\alpha/2) \sum_{i=1}^n |u_i|.
 \]
\end{lemma}
\noindent
Furthermore the following simple corollary of of the cosine rule will come in handy.
\begin{lemma}
\label{lem:smallanglediff}
 Let $z, z'$ be two complex numbers at an angle of at most $\pi/3$, then $|z-z'| \leq \max \{ |z|, |z'| \}$.
\end{lemma}
\begin{proof}
 Recall the cosine rule, for a triangle with sides $a$, $b$ and $c$; and angles $A$, $B$ and $C$ where side $a$ is not adjacent to angle $A$, then
 \[
 |a|^2 = |b|^2 + |c|^2 -2|b||c| \cos(A)\,,
 \]
 where $|a|$ is the length of side $a$.
 Now consider the triangle with vertices in $\CC$ at the origin, $z$ and $z'$.
 The sides have length $|z|, |z'|$ and $|z-z'|$ and the angle at the origin is the angle $\theta \leq \pi/3$ between $z$ and $z'$.
 As $\cos(x) \geq 1/2$ for $x \leq \pi/3$,
 \[
 |z-z'|^2 \leq |z|^2+|z'|^2-|z||z'| \leq \max\{|z|^2,|z'|^2\}\,.\qedhere
 \]
\end{proof}
 
To prove Lemma~\ref{lem:kcol} we need some definitions and an auxiliary lemma. 
We define rational functions (which depend on $k$ and $w$) in two variables $z_0,z$ and respectively $k-1$ variables $z_0,\ldots,z_{k-2}$ by
\begin{align*}
 R(z_0,z;w,k) &:= \frac{wz_0+(k-2)z+1}{z_0+(k-2)z+w}\,,\\
 R_k(z_0,z_1,\ldots,z_{k-2};w) &:= \frac{wz_0+z_1+\ldots+z_{k-2}+1}{z_0+z_1+\ldots+z_{z-2}+w}\,.
\end{align*}
Consider the cone
\[
C(\theta):=\{z=re^{i\vartheta}\mid r\geq 0 \text{ and } |\vartheta|\leq \theta\},
\]
and define for $d=0,\dotsc,\Delta$ and $c,\alpha> 0$, the region
\[
K(\theta,d,c,\alpha,\eps):=C(d\theta+\Delta-\eps)\cap \left\{z: \left(1+\frac{c}{\Delta}\right)^{d-\Delta}\left(1+\frac{\alpha}{\Delta}\right)^{d}\leq |z|\le\left(1+\frac{c}{\Delta}\right)^{\Delta-d}\left(1+\frac{\alpha}{\Delta}\right)^{d}\right\}\,.
\]
%
%
%
\begin{lemma}
\label{lem:forward invariant}
Let $\Delta\in \mathbb{N}_{\geq 3}$ and let $k\in \mathbb{N}_{\geq 3}$. Define $c=\log k-1$ and $\alpha=\log k^{1/2}-1$. Then there exists $0<\eps<\theta<\pi/(3\Delta)$ and $\eta=\omega(\frac{1}{\Delta})$
such that for each $d=0,\ldots,\Delta$, and any $z_0,\ldots,z_{k-2}\in K_d:=K(\theta,d,c,\alpha,\eps)$ such that for each $i,j$, $z_i/z_j\in K_d$ and any  $w\in \mathcal{N}([1,1+c/\Delta],\eta)$ the ratio $R=R_k(z_0,z_1,\ldots,z_{k-2};w)$ satisfies
\begin{align}
(1+\alpha/\Delta)^{-1}<|R|&< 1+\alpha/\Delta\quad \text{ and }\quad  |\arg(R)|< \theta.\label{eq:ratio bound arg}
\end{align}
In particular the values $\theta = 1/(5 \Delta)$, $\eps = \theta / [100 \log k]$ and $\eta = \min\{\Delta c / [800(\Delta + \alpha)^2], 1/[2400(\Delta+ \alpha)], c / [800 \Delta]\}$ suffice.
\end{lemma}
\noindent
We will prove this lemma in the next subsection, but we first utilize it to prove Lemma~\ref{lem:kcol}, which we restate here for convenience.
%
\zerofreeinduction*
%
\begin{proof}
The choice of constants is the same as in Lemma~\ref{lem:forward invariant} except that we need to choose $\eta$ small enough so that each $w\in \mathcal{N}([1,1+c/\Delta],\eta)$ has argument at most $\eps$. It thus suffices to take $\eta=\omega(\frac{1}{\Delta\log(k)})$.
%
We prove the lemma by induction on the number of free vertices of $G$.
For the base case, we have no free vertices and so every vertex is fixed.
Therefore $Z_L^W(G)$ is a product of non-zero terms, hence is non-zero, proving~\ref{itm:kcol,nonzero}.
Statement \ref{itm:kcol,leaf}\ref{itm:kcol,leaf,free} is vacuous as there are no free vertices.
Statement \ref{itm:kcol,leaf}\ref{itm:kcol,leaf,fixed} follows as the products $Z_{L'l}^{W'u}(G)$ and $Z_{L'l'}^{W'u}(G)$ differ in at most one term.
Thus their ratio is either $1$, $w$ or $w^{-1}$.
Similarly we deduce Statement~\ref{itm:kcol,nonleaf} (in which $d$ must be zero) from the fact that the products $Z_{L'l}^{W'u}(G)$ and $Z_{L'l'}^{W'u}(G)$ differ in at most $\Delta$ terms.
%
Now, we assume that Statements~\ref{itm:kcol,nonzero}, \ref{itm:kcol,leaf}, and~\ref{itm:kcol,nonleaf} hold for graphs with $r \ge 0$ free vertices. We prove the statements for $r+1$ free vertices. First, we shall prove~\ref{itm:kcol,nonzero}.
%
Suppose that $u$ is a free vertex.
Note that $Z_L^W(G) = \sum_{j=1}^k Z_{Lj}^{Wu}(G)$.
As each term in the sum on the right hand side of this expression has one fewer free vertex, we may apply induction to deduce that all of these terms are non-zero by \ref{itm:kcol,nonzero}.
Furthermore, by \ref{itm:kcol,nonleaf} each pair has angle at most $d\theta +(\Delta-d)\eps$ where $d$ is the number of free neighbours of $u$.
Lemma~\ref{lem:cone} tells us that the $Z_{Lj}^{Wu}$ all lie in a cone of angle at most $d\theta+(\Delta-d)\eps$ and
\begin{equation*}
 |Z_L^W(G)| = \bigg|\sum_{j=1}^k Z_{Lj}^{Wu}(G) \bigg| \geq \cos(d\theta/2+(\Delta-d)\eps/2)\sum_{j=1}^k |Z_{Lj}^{Wu}(G)| \neq 0\,.
\end{equation*}
%
Next, we shall prove~\ref{itm:kcol,leaf}\ref{itm:kcol,leaf,free} so consider the ratios,
\begin{align*}
 R_{j,l}(G) = \frac{Z_{L'j}^{W'u}(G)}{Z_{L'\ell}^{W'u}(G)}, & & R_{j,\ell}^v(G) = \frac{Z_{L'j}^{W'v}(G- u)}{Z_{L'\ell}^{W'v}(G - u)}.
\end{align*}
As $v$ is the unique neighbour of $u$ and is free, we may write, denoting $j^*$ for $\pi_{uv}(j)$ and $\ell^*$ for $\pi_{uv}(\ell)$,
\[
 R_{j,l}(G) = \frac{\sum_{i}Z^{Wuv}_{Lji}(G)}{\sum_{i}Z^{Wuv}_{L\ell i}(G)}
 =\frac{wZ^{Wv}_{Lj^*}(G-u)\sum_{i\notin\{ j^*,\ell^*\}}Z^{Wv}_{Li}(G-u)+Z^{Wv}_{L\ell^*}(G-u)}{Z^{Wv}_{Lj^*}(G-u)+\sum_{i\notin\{ j^*,\ell^*\}}Z^{Wv}_{Li}(G-u)+wZ^{Wv}_{L\ell^*}(G-u)}.
 \]
 Dividing both the numerator and denominator by $Z^{Wv}_{L\ell^*}(G-u)$ (which by induction is nonzero) we obtain,
 \begin{equation}
 \dfrac{w R_{j^*,\ell^*}^v(G) + \sum_{i \neq j^*, \ell^*} R_{i,\ell^*}^v(G) +1}{R_{j^*,\ell^*}^v(G) +\sum_{i \neq j^*,\ell^*} R_{i,l}^v(G)+w} = R_k(R_{j^*,\ell^*}^v(G), R_{1,\ell^*}^v(G), \dotsc, R_{k,\ell^*}^v(G);w)\,.
 \label{eq:recurrence}
\end{equation}
Where the function $R_k$ in~(\ref{eq:recurrence}) takes as arguments all $R^v_{i,\ell^*}(G)$ for $i \neq \ell^*$ precisely once (and so takes precisely $k-1$ arguments as expected.)
%
Suppose that $v$ has $d$ free neighbours that are not $u$.
Since $G -u$ has one fewer free vertex than $G$, we may apply the inductive hypothesis.
By \ref{itm:kcol,nonleaf} we find that for any $i\neq \ell^*$, we have $R^v_{i,\ell}(G)\in K_d$.
However, we also have that for any $i,j\neq \ell^*$, that 
\[
\frac{R^v_{i,\ell^*}(G)}{R^v_{j,\ell^*}(G)}=\frac{Z_{L'i}^{W'v}(G- u)}{Z_{L'j}^{W'v}(G- u)}=R^v_{i,j}(G)\in K_d.
\]
To prove~\ref{itm:kcol,leaf}\ref{itm:kcol,leaf,free}\ref{itm:kcol,leaf,free,angle}, observe that the angle between $Z_{L'j}^{W'u}$ and $Z_{L'l}^{W'u}$ is precisely the angle of $R_{j,l}(G)$ from the real axis in $\CC$ and so is bounded by the absolute value of the argument of $R_{j,l}(G)$, which by Lemma~\ref{lem:forward invariant} bounded by $\theta$ as desired.
Statement~\ref{itm:kcol,leaf}\ref{itm:kcol,leaf,free}\ref{itm:kcol,leaf,free,length} also follows immediately from Lemma~\ref{lem:forward invariant}.
%
For the proof of~\ref{itm:kcol,leaf}\ref{itm:kcol,leaf,fixed}, we note that as $v$ is fixed, then 
$$
Z_{L'j}^{W'u}(G) \in \{ w^{-1} Z_{L'l}^{W'u}(G),Z_{L'l}^{W'u}(G),wZ_{L'l}^{W'u}(G) \}
$$
from which both~\ref{itm:kcol,leaf,fixed,angle} and~\ref{itm:kcol,leaf,fixed,length} follow.
%
Finally, we prove \ref{itm:kcol,nonleaf}.
To do so we consider the graph $G \star u$ which is formed as follows.
Let $v_1, \ldots, v_r$ be the neighbours of $u$ ordered arbitrarily.
Let $u_1, \ldots, u_r$ be $r$ new vertices which will be copies of $u$.
Then $G \star u$ is the graph obtained by deleting $u$ and its incident edges, adding the vertices $u_1, \ldots, u_r$ and edges $u_1v_1, \ldots, u_rv_r$.
Furthermore, $G \star u$ inherits any colouring of $G$ and if $u$ is coloured, all of the new vertices inherit this colour.
Note that if $u$ is coloured, then the graph $G \star u$ has the same partition function as $G$.
Also, in this case $G \star u$ has the same number of free vertices as $G$.
This allows us to prove \ref{itm:kcol,nonleaf} from \ref{itm:kcol,leaf} by changing the colour of one copy of $u$ at a time.
That is,
\begin{equation}
 \frac{Z_{L'j}^{W'u}(G)}{Z_{L'l}^{W'u}(G)} = \frac{Z_{L'j\ldots j}^{W'u_1\ldots u_r}(G\star u )}{Z_{L'l\ldots l}^{W'u_1\ldots u_r}(G\star u )} = \prod_{i=1}^r \frac{Z_{L'j\ldots j l \ldots l}^{W'u_1\ldots u_{i-1}u_i \ldots u_r}(G\star u )}{Z_{L'j\ldots j l \ldots l}^{W'u_1\ldots u_i u_{i+1} \ldots u_r}(G\star u )} \label{eq:telescoping}
\end{equation}
By \ref{itm:kcol,leaf} each of the terms in the product in~(\ref{eq:telescoping}) has angle at most $\theta$ and absolute value at most $1 +\alpha/\Delta$ (if $u_i$ is free) or angle at most $\eps$ and absolute value at most $1 +c/\Delta$ (if $u_i$ is fixed).
As $u$ has $d$ free neighbours and at most $\Delta -d$ fixed neighbours, this allows us to conclude~\ref{itm:kcol,nonleaf}\ref{itm:kcol,nonleaf,angle} and \ref{itm:kcol,nonleaf}\ref{itm:kcol,nonleaf,length}, completing the induction.
\end{proof}
%
\subsection{Proof of Lemma~\ref{lem:forward invariant}}
\label{sec:boundcalpha}
%
We will require a technical lemma which concerns the real and imaginary parts of the ratios $R(z_1,z_2;w,k)$.
%
\begin{lemma}
\label{lem:reimparts}
 Let $z_1, z_2 \in \CC$ be defined as $z_1 = xe^{i\thx}$, $z_2 = ye^{i\thy}$ with $x,y \in \RR^+$ and $\thx,\thy \in [0,2\pi)$ and suppose $w \in[1, 1 + \frac{c}{\Delta}]$ is real.
 Then, the real and imaginary parts of $R(z_1,z_2;w,k)$ are as follows where $N$ is a nonzero constant,
 \begin{align}
 \label{eq:realR} \Re(R(z_1,z_2;w,k)) &= N(wx^2+(w+1)(k-2)xy\cos(\thx-\thy) +(k-2)^2y^2 \\ &\qquad+(w^2+1)x\cos(\thx)+(w+1)(k-2)y\cos(\thy) +w), \nonumber\\
 \label{eq:imagR} \Im(R(z_1,z_2;w,k)) &= N (w-1)((k-2)xy\sin(\thx-\thy) +(1+w)x\sin(\thx) +(k-2)y\sin(\thy)).
 \end{align}
\end{lemma}
\begin{remark} 
\label{rem:reimparts}
%
Set $\theta =\max(|\thx|,|\thy|, |\thx-\thy|)$ and assume $|\theta| \le 1$. Then as $|\sin t| \leq |t|$ and $\cos t \geq 1 - t^2/2$ for all $t$, and using $w \geq 1$ we obtain the following bounds:
 \begin{align*}
 \Re(R(z_1,z_2;w,k))  
 &\ge N(1 - \theta^2/2)(wx^2 + (w+1)(k-2)xy +(k-2)^2y^2  \\
&\qquad +(w^2+1)x +(w+1)(k-2)y +w) \\
 &\ge N(1 - \theta^2/2)(x+(k-2)y+w)(wx+(k-2)y+1);\\
 \text{and}\\
  \Im(R(z_1,z_2;w,k)) & \le N (w-1)((k-2)xy|\thx-\thy| +(1+w)x|\thx| +(k-2)y|\thy|).
 \end{align*} 
Hence
\begin{align}
 \bigg|\frac{\Im(R(z_1,z_2;w,k))}{\Re(R(z_1,z_2;w,k))}\bigg| \leq
 \frac{(w-1)\left((k-2)xy|\thx-\thy|+(1+w)x|\thx|+(k-2)y|\thy|\right)}
 {(1 - \frac{\theta^2}{2})(x+(k-2)y+w)(wx+(k-2)y+1)}. 
\label{eq:imreratio}
 \end{align}
\end{remark}
%
\begin{proof}
 We may write $z_1 = x\cos(\thx) +ix\sin(\thx)$ and $z_2 = y\cos(\thy) +iy\sin(\thy)$. Hence,
 \begin{align}
  \nonumber R(z_1,z_2;w,k) & = \frac{w(x\cos(\thx) +ix\sin(\thx)) +(k-2)(y\cos(\thy) +iy\sin(\thy)) +1}{x\cos(\thx) +ix\sin(\thx) +(k-2)(y\cos(\thy) +iy\sin(\thy)) +w} \\
  & = \frac{wx\cos(\thx) +(k-2)y\cos(\thy) +1 + i(wx\sin(\thx) + (k-2)y\sin(\thy))}{x\cos(\thx) +(k-2)y\cos(\thy) +w + i(x\sin(\thx) + (k-2)y\sin(\thy))} \label{eq:R(z1z2wk)1}
 \end{align}
Rationalising the denominator in~(\ref{eq:R(z1z2wk)1}), we obtain
\begin{align}
 \nonumber R(z_1,z_2;w,k) = & N^{-1}\left(wx\cos(\thx) +(k-2)y\cos(\thy) +1 + i(wx\sin(\thx) + (k-2)y\sin(\thy))\right)\\
 & \times \left(x\cos(\thx) +(k-2)y\cos(\thy) +w - i(x\sin(\thx) + (k-2)y\sin(\thy))\right) \label{eq:R(z1z2wk)2}
\end{align}
where $N = |x\cos(\thx) +(k-2)y\cos(\thy) +w + i(x\sin(\thx) + (k-2)y\sin(\thy))|^2$.
Expanding the expression in~(\ref{eq:R(z1z2wk)2}), the real and imaginary parts are given by the following expressions in which we write $c_x$ for $\cos(\thx)$ and similarly define $c_y,s_x$ and $s_y$ to simplify notation.
 \begin{align*}
  \Re(R(z_1,z_2;w,k)) &= N^{-1}(wx^2c_x^2+(w+1)(k-2)xyc_xc_y+(k-2)^2c_y^2 \\
  &\qquad + wx^2s_x^2 + (w+1)(k-2)xys_xs_y+(k-2)^2s_y^2\\
  &\qquad +(w^2+1)xc_x+(w+1)(k-2)yc_y+w)\\
  \Im(R(z_1,z_2;w,k)) &= N^{-1}((k-2)xy(c_xs_y+ws_xc_y) - (k-2)xy(wc_xs_y+s_xc_y)\\
  &\qquad + (w^2-1)xs_x+(w-1)(k-2)ys_y )
 \end{align*}
Combining these expressions with the trigonometric identities
\begin{gather*}
 \cos^2(\vartheta)+\sin^2(\vartheta) = 1 \\
 \sin(\alpha-\beta) = \sin(\alpha)\cos(\beta)-\sin(\beta)\cos(\beta) \\
 \cos(\alpha-\beta) = \cos(\alpha)\cos(\beta)+\sin(\alpha)\sin(\beta)
\end{gather*}
yields the expressions~(\ref{eq:realR}) and ~(\ref{eq:imagR}) as claimed.
%
By an application of the triangle law combined with an applications of the approximations, $|\sin\theta| \leq |\theta|$ and $\cos\theta \geq 1-\theta^2/2$, we obtain
\begin{align}
 |\Im(R(z_1,z_2;w,k))| &\leq N^{-1} (w-1)\left((k-2)xy|\thx-\thy|+(1+w)x|\thx|+(k-2)y|\thy|\right), \label{eq:Imbound} \\
 \Re(R(z_1,z_2;w,k)) &\geq N^{-1} \left((wx+(k-2)y+1)(x+(k-2)y+w) \right.\nonumber\\ 
&\qquad -\left. ((w+1)(k-2)(x+1)y+(w^2+1)x)\theta^2/2\right). \label{eq:Rebound1}
\end{align}
Dividing \eqref{eq:Imbound} by \eqref{eq:Rebound1}, noting that for $\theta$ small this is maximised when $w = 1+\frac{c}{\Delta}$ and regrouping some terms yields the bound \eqref{eq:imreratio}.
%
%
%
%
%
\end{proof}
\noindent
We can now give a proof of Lemma~\ref{lem:forward invariant}
%
\begin{proof}[Proof of Lemma~\ref{lem:forward invariant}]
We first prove a slightly stronger version of the lemma for $w'$ real.
That is, we will show that there exist a small constants $\kappa = c/100$ and $\kappa'=0.02$ such that
\begin{equation}\label{eq:real case}
(1+(\alpha-\kappa)/\Delta)^{-1}<|R|<1+(\alpha-\kappa)/\Delta \quad \text{ and}\quad |\arg(R)|< (1-\kappa')\theta.
\end{equation}
To do so we start by taking a constant $\delta$ small enough so that for all $k \geq 3$, $c$ and $\alpha$ satisfy the strict inequality 
%
\begin{equation}
 \frac{ce^c}{cos(\delta)(e^c+k-1)} < \alpha - \kappa; 
 \label{eq:calphabounds}
\end{equation}
for example $\delta = 1/2$ is sufficient.
%
%
%
Fix $d\in\{0,\ldots,\Delta\}$.
First we observe that we may assume that $|R|\geq 1$. 
Indeed, if $|R|<1$, then 
\[
1/R=\frac{z_0+\sum_{i=1}^{k-2}z_i+w}{wz_0+\sum_{i=1}^{k-2}z_i+1}=\frac{1+\sum_{i=1}^{k-2}z_i/z_0+w/z_0}{w+\sum_{i=1}^{k-2}z_i/z_0+1/z_0}
\]
and $|1/R|>1$.
Since for each $i,j\geq 0$, the pairs $z_i/z_0$ and $z_j/z_0$ also satisfy our assumptions this shows our claim.
We start by showing that $|R|$ is bounded by $1+(\alpha - \kappa)/\Delta$.
%
We observe that (setting $z = (z_1 + \cdots + z_{k-2})/k$)
\begin{equation}\label{eq:from 1}
|R| = |R(z_0,z;w,k)|=\bigg|1+\frac{(w-1)z_0+(1-w)}{z_0+\sum_{i=1}^{k-2}z_i+w}\bigg|\leq 1+\frac{\frac{c}{\Delta}|z_0-1|}{|z_0+\sum_{i=1}^{k-2}z_i+w|}.
\end{equation}
Lower bounding the denominator of~(\ref{eq:from 1}) may be done with an application of Barvinok's lemma.
For the numerator we apply Lemma~\ref{lem:smallanglediff} as the angle between $z_0$ and $1$ is certainly less than $\pi/3$. 
This allows us to deduce that
\[
 |R(z_0,z;w,k)| \leq 1 +\frac{\frac{c}{\Delta} \max\{|z_0|,1\}}{\cos(d\theta/2+(\Delta-d)\eps/2))(|z_0|+\sum_{i=1}^{k-2}|z_i|+1)}.
\]
We next observe that by symmetry we may assume that $|z_0|\leq 1$; otherwise we divide the numerator and the denominator by $z_0$.
To maximize the above quantity clearly one should take each $|z_i|$ as small as possible. 
So we take $|z_i| = (1+c/\Delta)^{d-\Delta}(1+\alpha/\Delta)^{-d} \geq (1+c/\Delta)^{-\Delta} \geq e^{-c}$ and noting that $\theta\leq \delta/\Delta$, we rearrange to deduce that 
\[
|R(z_0,z;w,k)| < 1+\frac{c / \Delta}{\cos(\delta)((k-1)e^{-c}+1)}< 1+(\alpha-\kappa)/\Delta
\]
by~(\ref{eq:calphabounds}). This proves the first bound in \eqref{eq:real case}.
%
%
%
%
%
For the other bound in \eqref{eq:real case}, recall that
 $z=\frac{1}{k-2}\sum_{i=1}^{k-2}z_j$ so that
 $R_k(z_0,z_1,\ldots,z_{k-2};w)=R(z_0,z;w,k)$. 
%
Note that $z\in C(d\theta+(\Delta-d)\eps)$ by convexity of the cone and so by Lemma~\ref{lem:cone} we have
\[
\cos(d\theta/2+(\Delta-d)\eps/2)(1+c/\Delta)^{d-\Delta}(1+\alpha/\Delta)^{-d}\leq |z|\leq (1+c/\Delta)^{\Delta-d}(1+\alpha/\Delta)^{d}.
\]
To prove the bound on the argument of $R(z_0,z;w,k)$ we use the inequality, $|\beta|\leq |\tan(\beta)|$. It therefore suffices to bound the ratio $\frac{|\Im R(z_0,z;w,k)|}{|\Re R(z_0,z;w,k)|} = \tan(\arg(R(z_0,z;w,k)))$, which by Lemma~\ref{lem:reimparts} and Remark~\ref{rem:reimparts} is bounded by
 \begin{equation}\label{eq:bound arg}
  \frac{(w-1)\left((k-2)|z_0z||\theta_0-\theta_z|+(1+w)|z_0 \theta_0|+(k-2)|z\theta_z|\right)}
 {(1 - \frac{\theta^2}{2})(|z_0|+(k-2)|z|+w)(w|z_0|+(k-2)|z|+1)}.  
\end{equation}
Now suppose we can prove that
\begin{equation}\label{eq:strict angle}
  \frac{\left((k-2)|z_0z||\theta_0-\theta_z|+(1+w)|z_0 \theta_0|+(k-2)|z\theta_z|\right)}
 {(|z_0|+(k-2)|z|+w)(w|z_0|+(k-2)|z|+1)}
  < \frac{\Delta \tau \theta}{c} 
\end{equation}
where $\tau  = 7/e^2 < 0.96$.
%
Then by choosing $\theta \leq 0.2$ and $\kappa' < 0.02$, we have that $\Delta \tau \theta /c < (1 - \frac{\theta^2}{2})(w-1)^{-1}(1 - \kappa')\theta$ (using $w < 1 + c/\Delta$). This together with \eqref{eq:strict angle}  proves that \eqref{eq:bound arg} is at most $\theta(1-\kappa')$ and hence $|\arg(R(z_0,z;w,k))|< \theta(1-\kappa')$, as desired.
%
We will now show that \eqref{eq:strict angle} holds.
So, first note that
\begin{align}
 \nonumber & \frac{\left((k-2)|z_0z||\theta_0-\theta_z|+(1+w)|z_0 \theta_0|+(k-2)|z\theta_z|\right)}
 {(|z_0|+(k-2)|z|+w)(w|z_0|+(k-2)|z|+1)} \\
 & \leq \frac{\left((k-2)|z_0z||\theta_0-\theta_z|+2|z_0\theta_0|+(k-2)|z\theta_z|\right)}{(|z_0|+(k-2)|z|+1)^2}, \label{eq:anglemax}
\end{align}
which can be observed by computing the derivative of the left hand side of~(\ref{eq:strict angle}) with respect to $w$ and noting it is strictly negative for $w \geq 1$.
Now, we maximize~(\ref{eq:anglemax}), so first we show that there is a maximum point where exactly two of $|\theta_0-\theta_z|,|\theta_0|,|\theta_z|$ are as large as possible and one is zero.
To see this, first note that clearly at least one of $|\theta_0 - \theta_z|, |\theta_0|, |\theta_z|$ must be as large as possible i.e. equal to $d\theta +(\Delta-d)\eps$.
In fact exactly two of these must be maximised as the maximization with respect to the $\theta$ terms only is of the form $f(\theta_0,\theta_z)=a|\theta_0 - \theta_z|+b|\theta_0| +c|\theta_z|$ for constants $a,b,c>0$.
So if $|\theta_0 - \theta_z| = d\theta +(\Delta-d)\eps$ for example, then if $b\geq c$ we may set $\theta_0=d\theta +(\Delta-d)\eps$, $\theta_z=0$ increasing $f(\theta_0,\theta_z)$.
Similar logic allows one to conclude that two of $|\theta_0 - \theta_z|, |\theta_0|, |\theta_z|$ are equal to $d\theta +(\Delta-d)\eps$ and one is $0$ in every other case.
%
This leaves us with three maximization problems over $R_d \subseteq \RR^2$ defined by 
\begin{align*}
 R_d = \{(x,y) |  (1+c/\Delta)^{d-\Delta}(1+\alpha/\Delta)^{-d} \leq & x \leq(1+c/\Delta)^{\Delta-d}(1+\alpha/\Delta)^{d} , \\
  \cos(d\theta/2+(\Delta-d)\eps/2)(1+c/\Delta)^{d-\Delta}(1+\alpha/\Delta)^{-d} \leq & y \leq (1+c/\Delta)^{\Delta-d}(1+\alpha/\Delta)^{d}, \\
  \cos(d\theta/2+(\Delta-d)\eps/2)(1+c/\Delta)^{d-\Delta}(1+\alpha/\Delta)^{-d} \leq & y/x\leq (1+c/\Delta)^{\Delta-d}(1+\alpha/\Delta)^{d}\}.
\end{align*}
We enlarge the region slightly obtaining the region $\widetilde{R}_d \subseteq \RR^2$ defined by
\begin{align*}
 \widetilde{R}_d & = \{(x,y) | \cos(\delta)exp(-(\frac{d}{\Delta} \alpha + (1-\frac{d}{\Delta})c)) \leq x, y, y/x \leq  exp(\frac{d}{\Delta} \alpha + (1-\frac{d}{\Delta})c)\cos(\delta)^{-1}\} \\
 & = \{(x,y) | \cos(\delta)e k^{\frac{d}{2\Delta}-1} \leq x, y, y/x \leq  k^{1-\frac{d}{2\Delta}}(e\cos(\delta))^{-1} \}
\end{align*}
The functions to maximise are,
\begin{align*}
 f_1(x,y) & =\frac{(k-2)(xy+y)}{(x+(k-2)y+1)^2}, \\
 f_2(x,y) & =\frac{(k-2)xy+2x}{(x+(k-2)y+1)^2}, \\
 f_3(x,y) & =\frac{2x+(k-2)y}{(x+(k-2)y+1)^2}.
\end{align*}
%
%
First we look at $f_1$, it has critical points along the line $x+1=(k-2)y$ where it attains its maximum value of $1/4$.
However, note that due to our choice of $c$ and $\alpha$, this line does not lie inside of $\widetilde{R_d}$, hence the maximum must be attained at a boundary point.
Furthermore both $f_2$ and $f_3$ have no critical points strictly inside the first quadrant, so again their maxima must be attained at a boundary point.
This allows us to reduce the problem to eighteen univariate maximization problems, each of which has maximum at most $3 e^{-1}k^{-\frac{d}{2\Delta}}$ over $\widetilde{R}_d$ (see Section~\ref{app:zfr} for details).
%
Thus~(\ref{eq:anglemax})  is upper bounded by  $(d\theta+(\Delta-d)\eps)3 e^{-1}k^{-\frac{d}{2\Delta}}$. As a function of $d$, this is maximised when $d = (\frac{2}{\log k} - \frac{\eps}{\theta - \eps})\Delta$, which  (if $d \geq 1$) gives an upper bound to \eqref{eq:anglemax} of 
\[
\frac{6e^{-2}\Delta(\theta - \eps)}{\log k} \exp\left(  \frac{\eps}{2(\theta - \eps)} \log k\right). 
\]
Thus \eqref{eq:strict angle} is satisfied provided $\frac{6}{7}(\theta - \eps)\exp( \frac{1}{2}\log k \frac{\eps}{\theta - \eps}) < \theta$. By taking $\eps = \theta x / \log k$ and assuming $\log k \geq 1$, the left hand side is bounded above by $\frac{6}{7} (1-x)\exp(x/2(1-x))\theta$ and this is at most $\theta$ (as required) by taking $x = 1/100$ as assumed in the statement of the lemma.   
%
%
%
%
%
%
%
%
%
If $d=0$, then as $f_1, f_2$ and $f_3$ are all bounded above by $1$, provided $\eps < \frac{\theta}{100\log(k)}$, the left hand side of~(\ref{eq:strict angle}) at most $\eps \Delta < \tau \theta \Delta/c$.
This completes the proof of~(\ref{eq:strict angle}) and hence of \eqref{eq:real case}.
%
We finally extend the proof to the case that $w\in \mathcal{N}([1,c/\Delta],\eta)$ 
for $\eta= 1/[800(\Delta + \alpha)^2]$ using continuity.
First observe that $R_k(w):=R_k(z_0, \ldots, z_{k-2};w)$ satisfies 
\[
R_k(w) = z_0 + \frac{(z_0 + (k-2)z + 1)(1 - z_0)}{z_0 + (k-2)z + w}.
\] 
Then 
\[
|R_k(w + \eta) - R_k(w)| 
= \left| \frac{[z_0 + (k-2)z + 1](1 - z_0)}{(z_0 + (k-2)z + w + \eta)(z_0 + (k-2)z + w)}\eta \right| 
%
\]
The numerator is upper bounded by  $[|z_0| + (k-2)|z| + 1](1 + |z_0|)|\eta|$, while the denominator is lower bounded by 
\[
\big[(|z_0| + (k-2)|z| + |w| - |\eta|\cos^{-1}(\Delta \theta/2))(|z_0| + (k-2)|z| + |w|)\cos(\Delta \theta/2)\big]^2
\]
where we use the fact that the angle between any two of $w, z_0, z$ is at most $\Delta \theta$ and so we can apply Barvinok's lemma. In the statement of the lemma, we assume $\Delta \theta \leq \pi / 3$ so $\cos(\Delta \theta) \geq 1/2$. Then using that $(x+a)/(x+b) \leq a/b$ for $x\geq 0$ and $a\geq b$ and using that $|w| \geq 1$ and that $\eta < 1/4$ (so that $|w| - |\eta|\cos^{-1}(\Delta \theta/2) > 1/2$), we have
\[
\frac{|z_0| + (k-2)|z| + 1}{(|z_0| + (k-2)|z| + |w| - |\eta| \cos^{-1}(\Delta \theta) ) \cos(\Delta \theta)}
\leq 4
 \:\:\:\text{and}\:\:\:
\frac{|z_0| +  1}{(|z_0| + (k-2)|z| + |w|)\cos(\Delta \theta)}
\leq 2.
\]
Combining the above inequalities we obtain $|R_k(w + \eta) - R_k(w)| \leq 8 \eta$. Recall $\eta \leq \min\{\Delta c / [800(\Delta + \alpha)^2], 1/[2400(\Delta+ \alpha)], c / [800 \Delta]\}$.
Then for $w\in \mathcal{N}([1,c/\Delta],\eta)$, we can write $w = w' + \eta$ with $w' \in [1,  \frac{c}{\Delta}]$ real. Writing $R = R(w)$, we have 
\[\left( 1 + \frac{\alpha}{\Delta} \right)^{-1}
\leq \left( 1 + \frac{\alpha - \kappa}{\Delta} \right)^{-1} - 8\eta
\leq |R(w')| - 8\eta
<|R| < |R(w)| + 8\eta 
\leq \left( 1 + \frac{\alpha - \kappa}{\Delta} \right) + 8\eta
\leq 
1 + \frac{\alpha}{\Delta},
\]
where the first and last inequalities follow by our choice of $\eta$.
%
It follows from simple geometry that if $|z_1 - z_2| \leq \mu$ for $z_1, z_2 \in \mathbb{C}$ and $\mu \in \mathbb{R}^+$ with $\mu < |z_1|$, then $|\arg(z_1) - \arg(z_2)| < \arcsin(\mu / |z_1|)$. Using this, and since $|R(w')|> (1 + \frac{\alpha}{\Delta})^{-1}$, we see that 
\[
\arg(R) < \arg(R(w')) + \arcsin\left( 8 \eta \left(1 + \frac{\alpha}{\Delta} \right) \right) = \theta(1 - \kappa')  + \arcsin\left( 8 \eta \left(1 + \frac{\alpha}{\Delta} \right) \right) < \theta;
\]
in order to check the last inequality holds, it is sufficient 
to check that $8 \eta (1 + \frac{\alpha}{\Delta}) < \sin(\kappa' \theta)$. 
Noting that $\sin x > x - x^3/6 > 5x/6$ for $x \in (0,1)$, 
it is sufficient that $ 8 \eta (1 + \frac{\alpha}{\Delta}) < 5\kappa'\theta / 6 $ (since $\kappa' \theta <1$ by our choice of $\kappa'$ and $\theta$) and this holds by our choice of $\eta$. This completes the proof of the lemma.
\end{proof}
%
\subsection{Maximization problems}
\label{app:zfr}
We look at the maximization problems coming from Section~\ref{sec:boundcalpha} and claim that each has an upper bound of at most $3 k^{-\frac{d}{2\Delta}}/e$.
We find eighteen of them, one for each of the three functions with either $x$, $y$, or $y/x$ fixed to one of the two corresponding boundary values.
This allows us to reduce to the univariate maximization problems detailed below.
To simplify the expressions we will let $r=k-2 $, $s = \cos(\delta)e k^{\frac{d}{2\Delta}-1}$ and $t=k^{1-\frac{d}{2\Delta}}(e\cos(\delta))^{-1}$.
%
\begin{center}
\begin{tabular}{ |c|c|c|c| }
\hline
& $f_1$ & $f_2$ & $f_3$ \\ \hline
  $x = s$ 
  & $p_1(y) =\frac{ry(1+s)}{(s+ry+1)^2}$
  & $p_2(y) = \frac{rsy+2s}{(s+ry+1)^2}$
  & $p_3(y) = \frac{2s+ry}{(s+ry+1)^2}$
  \\ \hline
  
  $x = t$ 
  & $p_4(y) = \frac{ry(1+t)}{(t+ry+1)^2}$ 
  & $p_5(y) = \frac{rty+2t}{(t+ry+1)^2}$
  & $p_6(y) = \frac{2t+ry}{(t+ry+1)^2}$
  \\ \hline
  
  $y = s$ 
  & $p_7(x) = \frac{rs(1+x)}{(x+rs+1)^2}$ 
  & $p_8(x) = \frac{rxs+2x}{(x+rs+1)^2}$
  & $p_9(x) = \frac{2x+rs}{(x+rs+1)^2}$
  \\ \hline
  
  $y = t$ 
  & $p_{10}(x) = \frac{rt(1+x)}{(x+rt+1)^2}$ 
  & $p_{11}(x) = \frac{rxt+2x}{(x+rt+1)^2}$
  & $p_{12}(x) = \frac{2x+rt}{(x+rt+1)^2}$
  \\ \hline
  
  $y/x = s$ 
  & $p_{13}(x) = \frac{rs(1+x^{-1})}{(x^{-1}+rs+1)^2}$ 
  & $p_{14}(x) = \frac{rs+2x^{-1}}{(x^{-1}+rs+1)^2}$ 
  & $p_{15}(x) = \frac{2x^{-1}+rx^{-1}s}{(x^{-1}+rs+1)^2}$ 
  \\ \hline
  
  $y/x = t$ 
  & $p_{16}(x) = \frac{rt(1+x^{-1})}{(x^{-1}+rt+1)^2}$ 
  & $p_{17}(x) = \frac{rt+2x^{-1}}{(x^{-1}+rt+1)^2}$ 
  & $p_{18}(x) = \frac{2x^{-1}+rx^{-1}t}{(x^{-1}+rt+1)^2}$  
  \\ \hline
\end{tabular}
\end{center}
%
To begin the maximization, first observe that under the map $x \mapsto x^{-1}$, each of the functions $p_j(x)$ is the same as some function $p_l(x)$ for some $13 \leq j \leq 18$ and $7 \leq l \leq 12$.
Furthermore, $y=s$ yields the bounds $s \leq x \leq 1$ and $y/x = s$ gives $1 \leq x \leq t$.
Similarly we may compare $y=t$ and $y/x = t$.
Thus the ranges for $x$ are identical after inverting $x$.
Hence we may ignore $p_{13}$ through $p_{18}$ leaving us with $12$ problems.
%
Next, consider $p_{10}$, $p_{11}$ and $p_{12}$, each of which can be bounded above by
$$
\frac{2rtx}{(x+rt+1)^2} \leq \frac{2rtx}{r^2t^2} \leq \frac{2}{r}
$$
where the final inequality follows as $x \leq t$.
%
Similarly, we can bound $p_4$, $p_5$ and $p_6$.
As it must be the case that $y \geq 1$, the numerator of each is bounded above by $2try$.
Thus an upper bound for all three is $2t/ry$.
Furthermore, $r \geq \frac{2k}{3\cos(\delta)}$ provided $k\geq 7$ and $\delta$ small enough. So we are left with an upper bound of $3 k^{-\frac{d}{2\Delta}}/e$.
%
The remaining problems are similar.
The numerators may all be bounded above by $rs(1+x) \leq 2rs$ (or for $p_1, p_2$ and $p_3$ by $2ry$.)
The denominators are all bounded from below by $r^2s^2$ and $r^2y^2$ respectively.
Thus all six of these are upper bounded by $2/rs$ which is at most $3 k^{-\frac{d}{2\Delta}}/e$.
%
Hence an upper bound on all of the problems $p_1$ through $p_{18}$ is $3 k^{-\frac{d}{2\Delta}}/e$ as claimed.
%
\subsection{Improvements for small \texorpdfstring{$k$}{k}}
When $k$ is small, then the parameter $c=\log(k)-1$ is also very small.
In fact we do not obtain a better constant than what is known for the Ising model until $k \geq 21$.
However it is possible to do better, we can choose different values for $\alpha$ and $c$ which work better in these cases.
In this section we will show how to derive such values.
%
First, we note that we may do the the analysis in an identical way until we find ourselves with the maximization problems $f_1, f_2$ and $f_3$.
Now we maximise these more carefully than in Section~\ref{app:zfr}.
First, for $f_1$ we apply AM-GM to the denominator to deduce that $f_1(x,y) \leq \frac{1}{4}$ for any $x,y$.
This allows us to take any $c<4$ and as $k$ is small this is all we need and so we may ignore this constraint.
This leaves us to maximise $f_2$ and $f_3$.
A similar argument to the one in the proof of Lemma~\ref{lem:forward invariant} allows us to deduce that the maxima are on the boundary of $R_d$ and hence we need only consider the boundary of $\widetilde{R_d}$.
%
Now, we proceed as in Section~\ref{app:zfr} with different choices of $s$ and $t$ where this time we will take $t=e^{d/\Delta \alpha + (1-d/\Delta)c}$ and $s=t^{-1}$.
We start with $12$ maximization problems which we reduce to $8$ by symmetry as before.
Furthermore, $f_2>f_3$ if and only if $x>1$ which allows us to halve the number of problems left to consider leaving us with $4$ problems.
More precisely, we are left with $p_3, p_5, p_9$ and $p_{11}$.
All of these are of the form $f(x)=(a_1 x+a_2)(x+a_3)^{-2}$ which has a maximum at $x=a_3-2a_2/a_1$.
See the following table for the maximization of these $4$ functions.
%
\begin{center}
\begin{tabular}{ |c|c|c|c|c|c| }
\hline
Function & $a_1$ & $a_2$ & $a_3$ & $x^*$ & $f(x^*)$ 
\\ \hline
$p_3$ & $\frac{1}{k-2}$ & $\frac{2s}{(k-2)^2}$ & $\frac{s+1}{k-2}$ & $\frac{1-3s}{k-2}$ & $\frac{1}{4(1-s)}$
\\ \hline
$p_5$ & $\frac{t}{k-2}$ & $\frac{2t}{(k-2)^2}$ & $\frac{t+1}{k-2}$ & $\frac{t-3}{k-2}\leq 1$ & $\frac{kt}{(t+k-1)^2}$
\\ \hline
$p_9$ & $2$ & $(k-2)s$ & $1+(k-2)s$ & $1$ & $\frac{1}{(2+(k-2)s)}$
\\ \hline
$p_{11}$ & $2+(k-2)t$ & $0$ & $1+(k-2)t$ & $1+(k-2)t>t$ & $\frac{(k-2)t^2+2t}{((k-1)t+1)^2}$
\\ \hline
\end{tabular}
\end{center}
Note that in the cases of $p_5$ and $p_{11}$ the maximum value $x_*$ is outside the domain which we are maximising over and thus we maximise at the endpoints of the domain instead.
%
Now, recall that the maximum values obtained above must also satisfy~(\ref{eq:calphabounds}).
Also, when $s=e^{-\alpha}$ it must be the case that $(2+(k-2)e^{-\alpha})^{-1} <c^{-1}$ (from $p_9$).
Combining these after rearrangement yields the inequity
\begin{equation}
 \frac{ce^c}{e^c+k-1} \leq \alpha \leq \log\bigg(\frac{k-2}{c-2}\bigg)
 \label{eq:howtofindca}
\end{equation}
We may solve this inequality computationally for $c$, and deduce that there is a choice of $\alpha, c$ which satisfies~(\ref{eq:howtofindca}) provided that $c \leq c_k$ for some $c_k$ which can be found in the following table.
The corresponding value of $\alpha_k$ is also provided.
We give both $c_k$ and $\alpha_k$ rounded to three decimal places.
%
\begin{center}
 \begin{tabular}{c||c|c|c|c|c|c|c|c|c|c}  $k$ & $3$ & $4$ & $5$ & $6$ & $7$ & $8$ & $9$ & $10$ & $11$ & $12$ \\
 \hline
  $\alpha_k$ & $1.767$ & $1.803$ & $1.849$ & $1.896$ & $1.944$ & $1.990$ & $2.034$ & $2.076$ & $2.116$ & $2.154$ \\  \hline
  $c_k$ & $2.171$ & $2.330$ & $2.472$ & $2.600$ & $2.716$ & $2.820$ & $2.916$ & $3.003$ & $3.084$ & $3.160$ \\
 \end{tabular}
\end{center}
%
Now, we check that these are indeed the maximum values. To do this, we first note that we have $p_3 \leq 1/4$ and applying AM-GM to the denominator of the maximum for $p_5$ yields a result which is smaller than the values from we obtained for the maximum of $p_9$.
Finally, for $p_{11}$, the denominator is at least $(k-1)(k-2) t^2 +2t(k-1)$. Thus, after cancellations we are left with $p_{11} \leq 1/(k-1)$ which suffices for $k \geq 4$.
For $k=3$ we can easily check that $(t^2+2t)(2t+1)^{-2}$ is maximised when $t=1$ and hence is certainly at most $1/3 < 1/2.17$.
%
Recall when computing the maximum of $p_9$, we took $s$ as large as possible where one would expect that we should do the opposite to maximise $p_9$.
We now justify this choice.
So recall that we must ensure $d\theta p_9(x) \leq \Delta \theta/c$.
Furthermore, $s$ may be considered as a function of $d$ and as such is equal to $\exp(-d/\Delta \alpha - (1-d/\Delta)c)$.
Thus we must ensure that
$$
g(d)=\frac{dc/\Delta}{2+(k-2)s} \leq 1.
$$
Writing $\lambda$ for $d/\Delta$ gives the following function with domain $[0,1]$
$$
G(\lambda)=\frac{\lambda c}{2+(k-2)e^{-\lambda \alpha -(1-\lambda)c}}.
$$
Differentiating this with respect to $\lambda$, we see that either $c-\alpha<1$ and $G$ is increasing on $[0,1]$ or there is a maximum with $\lambda >1$ which is not inside the domain.
Thus, we maximise $G$ at one of its boundary points and it is easy to see that $\lambda=1$ is the maximum point rather than $\lambda =0$ where $G(\lambda)=0$.
%
%
%
